\section{Job lifecycle and scheduling}

\subsection{State model}

The persisted state machine is authoritative. State transitions use compare-and-swap
semantics and append an event to the run history. Figure~\ref{fig:states} omits
administrative states for readability.

\begin{figure}[htbp]
\centering
\begin{tikzpicture}[
  node distance=8mm and 9mm,
  state/.style={draw=Primary,rounded corners,fill=LightBlue,minimum width=21mm,minimum height=8mm,align=center,font=\small},
  terminal/.style={draw=DarkGray,rounded corners,fill=LightGray,minimum width=21mm,minimum height=8mm,align=center,font=\small},
  flow/.style={-{Latex[length=2mm]},draw=DarkGray,thick},
  every edge quotes/.style={font=\scriptsize,fill=white,inner sep=1pt}
]
  \node[state] (created) {created};
  \node[state,right=of created] (blocked) {blocked};
  \node[state,right=of blocked] (ready) {ready};
  \node[state,right=of ready] (leased) {leased};
  \node[state,below=of leased] (running) {running};
  \node[state,left=of running] (finalizing) {finalizing};
  \node[terminal,left=of finalizing] (success) {succeeded};
  \node[terminal,below=of finalizing] (failed) {failed};
  \node[terminal,below=of blocked] (cancelled) {cancelled};

  \draw[flow] (created) -- (blocked);
  \draw[flow] (blocked) -- node[above,font=\scriptsize]{needs met} (ready);
  \draw[flow] (ready) -- (leased);
  \draw[flow] (leased) -- (running);
  \draw[flow] (running) -- (finalizing);
  \draw[flow] (finalizing) -- (success);
  \draw[flow] (finalizing) -- (failed);
  \draw[flow] (blocked) -- (cancelled);
  \draw[flow] (ready) |- (cancelled);
  \draw[flow] (leased) to[bend left=34] node[above,font=\scriptsize]{lease expires} (ready);
\end{tikzpicture}
\caption{Simplified job state machine}
\label{fig:states}
\end{figure}

A job may also end as \code{timed\_out}, \code{skipped}, or \code{infrastructure\_failed}.
An infrastructure failure is distinct from a user-command failure and is eligible for
automatic retry. Terminal states are immutable except through a documented
administrator correction process that appends, rather than rewrites, history.

\subsection{Lease protocol}

\begin{enumerate}
  \item An attested runner opens an authenticated long-poll request containing pool,
        capabilities, current concurrency, image digest, and health status.
  \item The scheduler selects a ready job and atomically creates an attempt with a unique
        fencing token and a default 60-second lease.
  \item The runner acknowledges within 10 seconds and renews every 20 seconds. Renewal
        proves possession of runner identity and the current fencing token.
  \item Missing two renewal windows makes the lease suspect; expiry returns the job to
        ready or creates a retry according to policy.
  \item Finalization is accepted only for the current attempt and fencing token. Late
        results are retained as diagnostic events but cannot change job state.
\end{enumerate}

External deployment effects cannot generally be rolled back by fencing alone. A
deployment attempt therefore carries an idempotency key derived from environment, run,
job, attempt, and artifact digest. Deployment adapters \must enforce that key or expose
a reconciliation operation.

\subsection{Scheduling policy}

Eligibility is a hard filter over pool, platform, architecture, labels, trust tier,
region, network zone, image policy, and resource request. Selection among eligible jobs
uses weighted fair queuing across organizations and repositories, with aging to prevent
starvation. Protected deployment jobs may use a reserved capacity class but cannot
preempt a running job.

For tenant $i$, the scheduler computes normalized dominant resource share

\[
  S_i = \max_{r \in R}\left(\frac{u_{i,r}}{w_i C_r}\right),
\]

where $u_{i,r}$ is current use of resource $r$, $C_r$ is allocatable capacity, and
$w_i$ is the tenant weight. Among ready queues, the scheduler prefers the lowest
$S_i$ after applying priority and age bounds. Policy limits prevent a tenant from
raising priority above its authorized class.

\subsection{Concurrency and cancellation}

A concurrency group serializes runs or selected jobs. Acquisition and release are
transactional. If \code{cancel\_in\_progress} is true, the prior run receives a
cancellation request and a bounded grace period; the successor does not enter a
protected environment until the previous deployment has stopped or reconciled.
Cancellation is cooperative for 10 seconds and then forceful. Cleanup steps run with
separate, tightly scoped credentials and a maximum duration.

% -----------------------------------------------------------------------------
